\documentclass[jou,11pt,donotrepeattitle]{apa7}
% APA 7th edition. Compilar:  latexmk -pdf backends.tex
\usepackage[american]{babel}
\usepackage{csquotes}
\usepackage{amsmath,amssymb}
\usepackage{booktabs}
\usepackage{graphicx}

\title{Hardware Reproducibility Noise Decays with Training: Measuring Backend
Effects Early Overstates Them by an Order of Magnitude}
\shorttitle{Backend Noise Decays with Training}
\leftheader{Speranza}

\author{Maximiliano Speranza}
\affiliation{Independent Researcher, Argentina}

\authornote{\addORCIDlink{Maximiliano Speranza}{0009-0005-0413-8554}
This work was carried out independently, without institutional affiliation, supervision, or
funding. The author is an undergraduate student of the University Technical Degree in
Programming (\emph{Tecnicatura Universitaria en Programaci\'on}, TUP) at Universidad
Tecnol\'ogica Nacional (UTN), Argentina; the programme neither sponsored nor supervised this
study, which is reported here in a personal capacity. Correspondence concerning this article
should be addressed to Maximiliano Speranza. E-mail: maximiliano.speranza@gmail.com.
The author declares no competing interests and received no funding. Code, stored validation
histories, and the script that reproduces every number reported here are publicly available
(see Open Practices).}

\abstract{%
Retraining an identical model with an identical seed on a different compute backend yields a
different validation curve, and it is often assumed that this irreproducibility is of the same
order as the seed-to-seed variation that papers already report. We measure it directly. Eight
seeds of a linear-attention model were trained to 7{,}500 steps on a Tesla T4 GPU and retrained
on CPU with identical code, data, and seeds, yielding 120 paired comparisons. The mean absolute
backend difference falls from $0.0331$ over steps 500--2{,}500 to $0.0060$ over steps
3{,}000--7{,}500, a $5.5\times$ decay: trajectories start apart and converge. Relative to the
between-seed standard deviation \emph{at the same step}, the backend term is $1.30\times$ in the
transient but only $0.53\times$ in the late regime, where architecture decisions are actually
taken. No directional bias is detectable ($t(7) = -1.42$, $p = .198$; 64 of 120 signed
differences negative). The practical consequence is a measurement caveat rather than a
reproducibility alarm: because between-seed noise is itself non-stationary---its standard
deviation falls from $0.101$ at step 500 to $0.011$ at step 7{,}500, a factor of nine---any
backend comparison anchored to a single global noise estimate is off by a large factor. We first
made that error ourselves and reported an inflated ratio of $3.00\times$; the correction is
documented here because the failure mode is easy to reproduce and hard to notice.}

\keywords{reproducibility, floating-point non-determinism, hardware variability,
early stopping, experimental methodology}

\begin{document}
\maketitle

\section{The question}

Two laboratories run the same code, with the same seed, on different GPUs, and obtain different
numbers. This is uncontroversial: floating-point reductions are not associative, and kernel
selection differs across devices. What is not established is \emph{how large} the resulting
disagreement is relative to the variation that papers already quantify, namely the spread across
random seeds. If the backend term were comparable to or larger than the seed term, then the
customary ``mean $\pm$ SD over $N$ seeds'' would omit a term of the same order, and architecture
comparisons whose reported margins are of that size would not be reproducible across machines.

We tested this on a concrete system rather than in general, and the answer depends critically on
\emph{when} the measurement is taken.

\section{Method}

The subject is condition \texttt{delta} of a pre-registered capacity experiment: a small
linear-attention model (193{,}493 parameters) trained on a multi-query associative recall task,
with validation accuracy evaluated every 500 steps at sequence loads 96 and 128. Eight seeds had
already been trained to 10{,}000 steps on a Tesla T4 GPU under JAX, with
\texttt{jax\_default\_matmul\_precision=highest}, \texttt{--xla\_gpu\_deterministic\_ops=true}
and \texttt{TF\_CUDNN\_DETERMINISTIC=1}.

Each seed was then retrained from scratch on CPU (Intel i3-7100, four threads) with identical
code, identical data generation, and the identical seed, to 7{,}500 steps. This yields 15 paired
evaluation points per seed and $8 \times 15 = 120$ paired comparisons. The comparison is
\emph{within seed}: every CPU value is compared against the GPU value of the same seed at the
same step, so the seed contribution is differenced out by construction.

Two quantities are reported at each step $s$: the mean absolute backend difference
$\overline{|d_s|} = \frac{1}{8}\sum_i |a^{\text{CPU}}_{i,s} - a^{\text{GPU}}_{i,s}|$, and the
between-seed standard deviation of the GPU runs \emph{at that same step}, $\sigma_s$. Their ratio
answers the question of interest: how much noise does changing the backend add, in units of the
noise the field already reports?

\section{Results}

\subsection{The backend effect decays with training}

\begin{table}[htbp]
\caption{Backend difference and between-seed dispersion, by training step}
\label{tab:main}
\small
\begin{tabular}{rccc}
\toprule
Step & $\overline{|d_s|}$ & $\sigma_s$ & Ratio \\
\midrule
500   & 0.0353 & 0.0972 & 0.36 \\
1{,}000  & 0.0362 & 0.0702 & 0.52 \\
1{,}500  & 0.0442 & 0.0371 & 1.19 \\
2{,}000  & 0.0275 & 0.0092 & 2.98 \\
2{,}500  & 0.0224 & 0.0133 & 1.68 \\
3{,}000  & 0.0199 & 0.0112 & 1.78 \\
3{,}500  & 0.0070 & 0.0145 & 0.48 \\
4{,}000  & 0.0060 & 0.0188 & 0.32 \\
5{,}000  & 0.0049 & 0.0090 & 0.54 \\
6{,}000  & 0.0030 & 0.0118 & 0.25 \\
7{,}000  & 0.0039 & 0.0103 & 0.38 \\
7{,}500  & 0.0035 & 0.0107 & 0.32 \\
\bottomrule
\end{tabular}

\vspace{5pt}
{\small\raggedright \textit{Note.} $\overline{|d_s|}$ is the mean absolute CPU--GPU difference in
validation accuracy across eight seeds; $\sigma_s$ is the between-seed standard deviation of the
GPU runs at the same step. Rows for 4{,}500, 5{,}500 and 6{,}500 omitted for space; the full
table is in the repository.\par}
\end{table}

Grouping into regimes, the mean absolute difference is $0.0331$ over steps 500--2{,}500 and
$0.0060$ over steps 3{,}000--7{,}500, a decay of $5.5\times$. Expressed relative to between-seed
dispersion at the same step, the backend term averages $1.30\times$ in the transient and
$0.53\times$ in the late regime. Over all 120 comparisons the ratio is $0.81\times$.

The interpretation is that the two trajectories are not diverging but converging. Different
arithmetic on different devices separates the runs early, when the loss surface is being
traversed quickly and small perturbations are amplified; as training settles, both runs approach
the same region and the difference shrinks.

\subsection{No directional bias}

Of the 120 signed differences, 64 are negative. Averaging within seed and testing the eight
per-seed means against zero gives $t(7) = -1.42$, $p = .198$. There is no evidence that CPU
systematically underperforms or outperforms the GPU; what the backend contributes is dispersion,
not a shift. This conclusion was stable across sample sizes ($n = 3, 4, 8$) and across both
horizons, and it is the most robust finding of the experiment.

\subsection{A control that came for free}

One seed lacked a checkpoint and was retrained from zero on the same machine. It reproduced
$0.5182$ at step 500 against $0.51822$ in the original CPU run. Within a backend, training is
exactly reproducible; the differences measured here therefore arise from the change of device
and not from run-to-run variability.

\section{The error we made first, and why it is worth reporting}

Our first analysis compared the backend differences, which were measured over steps 500--2{,}500,
against a \emph{single} between-seed standard deviation of $0.01105$---the value at step 7{,}500.
That produced a ratio of $3.00\times$ and a worst case of $20.25\times$, and supported the
conclusion that backend choice dominates seed choice.

It is wrong, because between-seed dispersion is not a property of the experiment but a function
of the step:

\begin{center}
\small
\begin{tabular}{rccccccc}
\toprule
Step & 500 & 1{,}000 & 1{,}500 & 2{,}000 & 2{,}500 & 5{,}000 & 7{,}500 \\
\midrule
$\sigma_s$ & .101 & .073 & .038 & .010 & .014 & .009 & .011 \\
\bottomrule
\end{tabular}
\end{center}

A factor of nine between the transient and the late regime. Anchoring an early-regime numerator
to a late-regime denominator inflates the ratio; pairing each step against its own $\sigma_s$
brings $3.00\times$ down to $1.30\times$ over the same steps. The reported worst case was the
same artifact: the largest single difference, $0.224$, occurs at step 1{,}500 of one seed, where
$\sigma_s = 0.038$, not $0.011$.

We report this because the failure mode is easy to reproduce, invisible in the output, and would
have led us to publish the opposite of what the extended measurement shows. Any study that
summarizes non-stationary noise with one number is exposed to it.

\section{Limitations}

The result is specific to one architecture, one task, one GPU model and one CPU. It does not
separate ``CPU versus GPU'' from ``this CPU versus this GPU,'' and it says nothing about
mixed-precision training, distributed reductions, or larger models, where the mechanism may
differ in magnitude or in sign. The CPU runs were executed with the processor frequency capped
at 75\% for thermal reasons; this changes wall-clock time per step, not arithmetic. Finally,
$p = .198$ with eight seeds does not establish the absence of bias---it establishes that any
bias is small relative to dispersion.

\section{Implications}

The headline is a measurement caveat, not a reproducibility alarm. In the regime where
architecture comparisons are actually decided, changing backend contributed roughly half the
variation of changing seed, and about 30\% of the decision margin of the study that hosted this
experiment. Reporting ``mean $\pm$ SD over $N$ seeds'' is not meaningfully incomplete on that
account.

The transferable point is narrower and, we think, more useful: \emph{reproducibility studies must
declare the training regime in which they measured}. A backend comparison conducted at 2{,}500
steps and one conducted at 7{,}500 steps on the same system differ by a factor of five, and the
early number is the alarming one. Absent that declaration, a reader cannot tell whether a
reported irreproducibility is a property of the system or an artifact of when the ruler was
applied.

\section*{Open Practices}

The training code, the stored validation histories for all sixteen runs, the analysis script that
regenerates every number in this note, and the full step-by-step table are available at
\texttt{github.com/SperanzaMax/telar-ligamento} under \texttt{resultados/backend/}. The
non-stationarity of validation noise documented in Section~5 motivated a feature of the
\texttt{stoppower} package (\url{https://doi.org/10.5281/zenodo.21711767}), which requires the
user to declare the step range over which noise is estimated and warns when the constant-noise
assumption fails within a window. The pre-registered study that supplied the runs is reported
separately (\url{https://doi.org/10.5281/zenodo.21630279}).

\end{document}
